\section{Results}\label{results}

Results from the experiments conducted on Massachusetts and Cook County data are summarized in Table~\ref{tab:bpr_massachusetts} and Table~\ref{tab:bpr_cookcounty}, respectively. The best model(s) can vary depending on the chosen evaluation metric.

In both catchment areas, we see reasons to prefer our proposed BPR metric to alternatives when the goal is effectively prioritizing where to intervene. First, in Massachusetts, both the Gaussian Process (GP) and Bayesian Spatiotemporal model (BST) have top performance as assessed by MAE and RMSE. However, the BST has higher \%BPR than the GP (62.0\% compared to 58.2\%). Interventions guided by the BST model would have the potential to preemptively identify 18 additional fatal overdoses per year in the top 100 census tracts. Similarly, in Cook County the Gradient Boosted Trees model with SVI covariates has superior MAE and RMSE to the GLM model, yet has worse \%BPR (77.1\% versus 79.4\% for GLM). Interventions guided by the GLM model (preferred via the \%BPR metric) could reach 15 more fatal overdose events annually in Cook County.

We also observe that while complex models like BST do well in both catchment areas, so does the simple \emph{historical average} baseline and its weighted extension. In Massachusetts, \emph{historical average} delivers a BPR and MAE that fall within the uncertainty intervals of the best performing models. The weighted extension's ultimate \%BPR is so close to the top method (NPSpLag) that the difference amounts to identifying fewer than 2 additional overdose events annually. In Cook County, the \emph{historical average} baseline delivers competitive scores (within the intervals of the best models) as assessed by all three metrics (BPR, MAE, and RMSE); the best model (BST) here would reach less than 1 additional overdose event annually.

\begin{longtable}[]{@{}
  >{\raggedright\arraybackslash}p{(\linewidth - 8\tabcolsep) * \real{0.2308}}
  >{\raggedright\arraybackslash}p{(\linewidth - 8\tabcolsep) * \real{0.1923}}
  >{\raggedright\arraybackslash}p{(\linewidth - 8\tabcolsep) * \real{0.1923}}
  >{\raggedright\arraybackslash}p{(\linewidth - 8\tabcolsep) * \real{0.1923}}
  >{\raggedright\arraybackslash}p{(\linewidth - 8\tabcolsep) * \real{0.1923}}@{}}
\caption{Comparison of fatal opioid-related overdose prediction models trained on Massachusetts decedent data from 2010-2019, then evaluated on data from 2020 and 2021.}\label{tab:bpr_massachusetts}\\
\toprule\noalign{}
\begin{minipage}[b]{\linewidth}\raggedright
\end{minipage} & \multicolumn{4}{>{\centering\arraybackslash}p{(\linewidth - 8\tabcolsep) * \real{0.7692} + 6\tabcolsep}@{}}{%
\begin{minipage}[b]{\linewidth}\centering
Metric
\end{minipage}} \\
\begin{minipage}[b]{\linewidth}\raggedright
Model
\end{minipage} & \begin{minipage}[b]{\linewidth}\raggedright
Total overdoses identified

in top 100 tracts
\end{minipage} & \begin{minipage}[b]{\linewidth}\raggedright
\%BPR (K=100)
\end{minipage} & \begin{minipage}[b]{\linewidth}\raggedright
MAE
\end{minipage} & \begin{minipage}[b]{\linewidth}\raggedright
RMSE
\end{minipage} \\
\begin{minipage}[b]{\linewidth}\raggedright
All Zeros
\end{minipage} & \begin{minipage}[b]{\linewidth}\raggedleft
124.1
\end{minipage} & \begin{minipage}[b]{\linewidth}\raggedleft
\hl{25.1, (24.8-25.8)}
\end{minipage} & \begin{minipage}[b]{\linewidth}\raggedleft
\hl{1.24, (1.18-1.30)}
\end{minipage} & \begin{minipage}[b]{\linewidth}\raggedleft
\hl{1.92, (1.83-2.01)}
\end{minipage} \\
\begin{minipage}[b]{\linewidth}\raggedright
Last year
\end{minipage} & \begin{minipage}[b]{\linewidth}\raggedleft
254.5
\end{minipage} & \begin{minipage}[b]{\linewidth}\raggedleft
\hl{51.8, (49.5-54.4)}
\end{minipage} & \begin{minipage}[b]{\linewidth}\raggedleft
\hl{1.09, (1.04-1.14)}
\end{minipage} & \begin{minipage}[b]{\linewidth}\raggedleft
\hl{1.58, (1.50-1.65)}
\end{minipage} \\
\begin{minipage}[b]{\linewidth}\raggedright
Historical Average

(4 year)
\end{minipage} & \begin{minipage}[b]{\linewidth}\raggedleft
295.4
\end{minipage} & \begin{minipage}[b]{\linewidth}\raggedleft
\textbf{\hl{59.8, (56.8-62.9)}}
\end{minipage} & \begin{minipage}[b]{\linewidth}\raggedleft
\textbf{0.92, (0.90-0.94)}
\end{minipage} & \begin{minipage}[b]{\linewidth}\raggedleft
1.28, (1.24-1.33)
\end{minipage} \\
\begin{minipage}[b]{\linewidth}\raggedright
Weighted

Historical Average
\end{minipage} & \begin{minipage}[b]{\linewidth}\raggedleft
303.5
\end{minipage} & \begin{minipage}[b]{\linewidth}\raggedleft
\textbf{\hl{61.4, (56.1-66.9)}}
\end{minipage} & \begin{minipage}[b]{\linewidth}\raggedleft
\textbf{0.94, (0.91-0.98)}
\end{minipage} & \begin{minipage}[b]{\linewidth}\raggedleft
1.32, (1.27-1.38)
\end{minipage} \\
\begin{minipage}[b]{\linewidth}\raggedright
Poisson GLM
\end{minipage} & \begin{minipage}[b]{\linewidth}\raggedleft
304.7
\end{minipage} & \begin{minipage}[b]{\linewidth}\raggedleft
\textbf{61.6, (56.4-66.8)}
\end{minipage} & \begin{minipage}[b]{\linewidth}\raggedleft
\textbf{\hl{0.95, (0.91-0.98)}}
\end{minipage} & \begin{minipage}[b]{\linewidth}\raggedleft
\hl{1.32, (1.25-1.39)}
\end{minipage} \\
\begin{minipage}[b]{\linewidth}\raggedright
Poisson GLM +SVI
\end{minipage} & \begin{minipage}[b]{\linewidth}\raggedleft
301.3
\end{minipage} & \begin{minipage}[b]{\linewidth}\raggedleft
\textbf{61.0, (57.4-65.1)}
\end{minipage} & \begin{minipage}[b]{\linewidth}\raggedleft
\hl{1.10, (1.04-1.16)}
\end{minipage} & \begin{minipage}[b]{\linewidth}\raggedleft
\hl{1.38, (1.30-1.47)}
\end{minipage} \\
\begin{minipage}[b]{\linewidth}\raggedright
Gradient Boosted Trees +SVI
\end{minipage} & \begin{minipage}[b]{\linewidth}\raggedleft
293.6
\end{minipage} & \begin{minipage}[b]{\linewidth}\raggedleft
\textbf{59.5, (55.9-63.1)}
\end{minipage} & \begin{minipage}[b]{\linewidth}\raggedleft
\textbf{0.92, (0.89-0.96)}
\end{minipage} & \begin{minipage}[b]{\linewidth}\raggedleft
\textbf{1.24, (1.18-1.30)}
\end{minipage} \\
\begin{minipage}[b]{\linewidth}\raggedright
GP
\end{minipage} & \begin{minipage}[b]{\linewidth}\raggedleft
287.2
\end{minipage} & \begin{minipage}[b]{\linewidth}\raggedleft
58.2, (55.2-61.2)
\end{minipage} & \begin{minipage}[b]{\linewidth}\raggedleft
\textbf{0.93, (0.91-0.95)}
\end{minipage} & \begin{minipage}[b]{\linewidth}\raggedleft
\textbf{1.28, (1.23-1.34)}
\end{minipage} \\
\begin{minipage}[b]{\linewidth}\raggedright
CASTNet +SVI
\end{minipage} & \begin{minipage}[b]{\linewidth}\raggedleft
268.0
\end{minipage} & \begin{minipage}[b]{\linewidth}\raggedleft
54.4, (52.0-56.6)
\end{minipage} & \begin{minipage}[b]{\linewidth}\raggedleft
1.07, (0.99-1.16)
\end{minipage} & \begin{minipage}[b]{\linewidth}\raggedleft
1.47, (1.36-1.58)
\end{minipage} \\
\begin{minipage}[b]{\linewidth}\raggedright
BST: Bayesian spatiotemporal
\end{minipage} & \begin{minipage}[b]{\linewidth}\raggedleft
301.0
\end{minipage} & \begin{minipage}[b]{\linewidth}\raggedleft
\textbf{61.1, (58.6-63.3)}
\end{minipage} & \begin{minipage}[b]{\linewidth}\raggedleft
0.95, (0.93--0.96)
\end{minipage} & \begin{minipage}[b]{\linewidth}\raggedleft
1.26, (1.24-1.28)
\end{minipage} \\
\begin{minipage}[b]{\linewidth}\raggedright
BST + SVI
\end{minipage} & \begin{minipage}[b]{\linewidth}\raggedleft
305.4
\end{minipage} & \begin{minipage}[b]{\linewidth}\raggedleft
\textbf{62.0, (60.0-63.7)}
\end{minipage} & \begin{minipage}[b]{\linewidth}\raggedleft
\textbf{0.92, (0.91-0.94)}
\end{minipage} & \begin{minipage}[b]{\linewidth}\raggedleft
\textbf{1.23, (1.21-1.25)}
\end{minipage} \\
\begin{minipage}[b]{\linewidth}\raggedright
NBSpLag + SVI
\end{minipage} & \begin{minipage}[b]{\linewidth}\raggedleft
305.4
\end{minipage} & \begin{minipage}[b]{\linewidth}\raggedleft
\textbf{62.0, (60.3-64.0)}
\end{minipage} & \begin{minipage}[b]{\linewidth}\raggedleft
\textbf{0.92 (0.89-0.95)}
\end{minipage} & \begin{minipage}[b]{\linewidth}\raggedleft
1.31, (1.27-1.35)
\end{minipage} \\
\midrule\noalign{}
\endhead
\bottomrule\noalign{}
\endlastfoot
\end{longtable}


\begin{longtable}[]{@{}
  >{\raggedright\arraybackslash}p{(\linewidth - 8\tabcolsep) * \real{0.2308}}
  >{\raggedright\arraybackslash}p{(\linewidth - 8\tabcolsep) * \real{0.1875}}
  >{\raggedright\arraybackslash}p{(\linewidth - 8\tabcolsep) * \real{0.1939}}
  >{\raggedright\arraybackslash}p{(\linewidth - 8\tabcolsep) * \real{0.1955}}
  >{\raggedright\arraybackslash}p{(\linewidth - 8\tabcolsep) * \real{0.1923}}@{}}
\caption{Comparison of fatal opioid-related overdose prediction models trained on Cook County, Illinois decedent data, from 2015 to 2020, then evaluated on data from 2021 and 2022.}\label{tab:bpr_cookcounty}\\
\toprule\noalign{}
\begin{minipage}[b]{\linewidth}\raggedright
\end{minipage} & \multicolumn{4}{>{\centering\arraybackslash}p{(\linewidth - 8\tabcolsep) * \real{0.7692} + 6\tabcolsep}@{}}{%
\begin{minipage}[b]{\linewidth}\centering
Metric
\end{minipage}} \\
\begin{minipage}[b]{\linewidth}\raggedright
Model
\end{minipage} & \begin{minipage}[b]{\linewidth}\raggedright
Total overdoses

identified in top 100 tracts
\end{minipage} & \begin{minipage}[b]{\linewidth}\raggedright
\%BPR (K=100)
\end{minipage} & \begin{minipage}[b]{\linewidth}\raggedright
MAE
\end{minipage} & \begin{minipage}[b]{\linewidth}\raggedright
RMSE
\end{minipage} \\
\begin{minipage}[b]{\linewidth}\raggedright
All Zeros
\end{minipage} & \begin{minipage}[b]{\linewidth}\raggedleft
137.0
\end{minipage} & \begin{minipage}[b]{\linewidth}\raggedright
\hl{21.7, (21.1--22.4)}
\end{minipage} & \begin{minipage}[b]{\linewidth}\raggedright
1.37, (1.30--1.44)
\end{minipage} & \begin{minipage}[b]{\linewidth}\raggedright
2.46, (2.34--2.56)
\end{minipage} \\
\begin{minipage}[b]{\linewidth}\raggedright
Last year
\end{minipage} & \begin{minipage}[b]{\linewidth}\raggedleft
466.0
\end{minipage} & \begin{minipage}[b]{\linewidth}\raggedright
\hl{73.9, (70.7--77.0)}
\end{minipage} & \begin{minipage}[b]{\linewidth}\raggedright
\hl{1.07, (1.03--1.10)}
\end{minipage} & \begin{minipage}[b]{\linewidth}\raggedright
1.66, (1.59--1.71)
\end{minipage} \\
\begin{minipage}[b]{\linewidth}\raggedright
Historical Average (4 year)
\end{minipage} & \begin{minipage}[b]{\linewidth}\raggedleft
505.3
\end{minipage} & \begin{minipage}[b]{\linewidth}\raggedright
\textbf{\hl{80.1, (76.2--84.3)}}
\end{minipage} & \begin{minipage}[b]{\linewidth}\raggedright
\textbf{0.94, (0.90--0.97)}
\end{minipage} & \begin{minipage}[b]{\linewidth}\raggedright
\textbf{1.44, (1.38--1.50)}
\end{minipage} \\
\begin{minipage}[b]{\linewidth}\raggedright
Weighted

Historical Average
\end{minipage} & \begin{minipage}[b]{\linewidth}\raggedleft
495.4
\end{minipage} & \begin{minipage}[b]{\linewidth}\raggedright
\textbf{\hl{78.6, (75.7--81.8)}}
\end{minipage} & \begin{minipage}[b]{\linewidth}\raggedright
\hl{1.15, (1.11--1.19)}
\end{minipage} & \begin{minipage}[b]{\linewidth}\raggedright
1.87, (1.76--1.96)
\end{minipage} \\
\begin{minipage}[b]{\linewidth}\raggedright
Poisson GLM
\end{minipage} & \begin{minipage}[b]{\linewidth}\raggedleft
496.8
\end{minipage} & \begin{minipage}[b]{\linewidth}\raggedright
\textbf{78.8, (75.9--82.3)}
\end{minipage} & \begin{minipage}[b]{\linewidth}\raggedright
\hl{1.22, (1.08--1.29)}
\end{minipage} & \begin{minipage}[b]{\linewidth}\raggedright
\hl{4.64, (1.65--5.92)}
\end{minipage} \\
\begin{minipage}[b]{\linewidth}\raggedright
Poisson GLM + SVI
\end{minipage} & \begin{minipage}[b]{\linewidth}\raggedleft
500.7
\end{minipage} & \begin{minipage}[b]{\linewidth}\raggedright
\textbf{79.4, (76.0--82.8)}
\end{minipage} & \begin{minipage}[b]{\linewidth}\raggedright
1.16, (1.11--1.20)
\end{minipage} & \begin{minipage}[b]{\linewidth}\raggedright
\hl{1.86, (1.77--1.94)}
\end{minipage} \\
\begin{minipage}[b]{\linewidth}\raggedright
Gradient Boosted Trees +SVI
\end{minipage} & \begin{minipage}[b]{\linewidth}\raggedleft
485.7
\end{minipage} & \begin{minipage}[b]{\linewidth}\raggedright
\textbf{77.1, (72.2--81.9)}
\end{minipage} & \begin{minipage}[b]{\linewidth}\raggedright
\hl{0.99, (0.95--1.05)}
\end{minipage} & \begin{minipage}[b]{\linewidth}\raggedright
\hl{1.55, (1.42--1.68)}
\end{minipage} \\
\begin{minipage}[b]{\linewidth}\raggedright
GP
\end{minipage} & \begin{minipage}[b]{\linewidth}\raggedleft
477.2
\end{minipage} & \begin{minipage}[b]{\linewidth}\raggedright
\textbf{75.7, (69.9--80.5)}
\end{minipage} & \begin{minipage}[b]{\linewidth}\raggedright
\hl{1.03, (0.97--1.09)}
\end{minipage} & \begin{minipage}[b]{\linewidth}\raggedright
1.63, (1.50--1.74)
\end{minipage} \\
\begin{minipage}[b]{\linewidth}\raggedright
CASTNet +SVI
\end{minipage} & \begin{minipage}[b]{\linewidth}\raggedleft
472.6
\end{minipage} & \begin{minipage}[b]{\linewidth}\raggedright
75.2, (73.2--76.8)
\end{minipage} & \begin{minipage}[b]{\linewidth}\raggedright
1.01, (0.98-1.04)
\end{minipage} & \begin{minipage}[b]{\linewidth}\raggedright
1.53, (1.39-1.67)
\end{minipage} \\
\begin{minipage}[b]{\linewidth}\raggedright
BST: Bayesian spatiotemporal
\end{minipage} & \begin{minipage}[b]{\linewidth}\raggedleft
505.9
\end{minipage} & \begin{minipage}[b]{\linewidth}\raggedright
\textbf{80.5, (78.9--82.1)}
\end{minipage} & \begin{minipage}[b]{\linewidth}\raggedright
1.00, (0.98--1.01)
\end{minipage} & \begin{minipage}[b]{\linewidth}\raggedright
\textbf{1.40, (1.36--1.43)}
\end{minipage} \\
\begin{minipage}[b]{\linewidth}\raggedright
BST + SVI
\end{minipage} & \begin{minipage}[b]{\linewidth}\raggedleft
504.1
\end{minipage} & \begin{minipage}[b]{\linewidth}\raggedright
\textbf{80.2, (78.7--82.0)}
\end{minipage} & \begin{minipage}[b]{\linewidth}\raggedright
0.97, (0.95--0.98)
\end{minipage} & \begin{minipage}[b]{\linewidth}\raggedright
1.47, (1.43--1.51)
\end{minipage} \\
\begin{minipage}[b]{\linewidth}\raggedright
NBSpLag + SVI
\end{minipage} & \begin{minipage}[b]{\linewidth}\raggedleft
502.2
\end{minipage} & \begin{minipage}[b]{\linewidth}\raggedright
\textbf{\hl{79.9, (78.1--81.6)}}
\end{minipage} & \begin{minipage}[b]{\linewidth}\raggedright
\textbf{\hl{0.96, (0.93-0.98)}}
\end{minipage} & \begin{minipage}[b]{\linewidth}\raggedright
\textbf{\hl{1.42, (1.37-1.48)}}
\end{minipage} \\
\midrule\noalign{}
\endhead
\bottomrule\noalign{}
\endlastfoot
\end{longtable}

